\chapter{Introdução}
\label{chap:introducao}

O mercado de criptoativos consolidou-se, ao longo da última década, como uma
das classes de ativos mais dinâmicas e inovadoras do sistema financeiro
global. Entre esses ativos, o Bitcoin ocupa posição central em termos de
capitalização, liquidez e profundidade de mercado, atuando como ativo de
\textit{benchmark} para o desenvolvimento de derivativos, índices e estratégias
quantitativas.

Desde a proposta original de \citet{nakamoto2008bitcoin}, o Bitcoin evoluiu de
um experimento de moeda digital descentralizada para um ativo financeiro com um
ecossistema robusto de derivativos, destacando-se a negociação de contratos
futuros e de opções listadas em plataformas especializadas, como a Deribit. A
expansão desse mercado tornou possível observar, em alta frequência, a formação
das expectativas de volatilidade por parte de investidores institucionais e
\textit{market makers}. A volatilidade implícita extraída dos preços de opções
sintetiza tanto a visão agregada do mercado sobre a incerteza futura quanto a
demanda por proteção contra movimentos adversos do preço do ativo subjacente.

Nesse contexto, o \textit{Volatility Risk Premium} (VRP) --- usualmente definido
como a diferença entre a volatilidade implícita e a volatilidade realizada
observada ex post --- assume papel fundamental na precificação de risco. Em
mercados tradicionais, como ações e índices, a evidência empírica documenta que
o VRP tende a ser sistematicamente positivo, refletindo o prêmio pago por agentes
dispostos a adquirir proteção contra quedas abruptas de preços e contra choques
súbitos de volatilidade \citep{bollerslev2011VRP, bali2009volatilitypremium,
pan2006jumpRisk}. Além disso, diversos estudos apontam que o VRP carrega
informação relevante para a previsão de retornos, a identificação de regimes de
mercado e a avaliação de condições de estresse financeiro.

Apesar da rápida evolução do mercado de derivativos de criptoativos, a dinâmica
do VRP em Bitcoin permanece relativamente pouco investigada na literatura
acadêmica. Pesquisas recentes começam a documentar a existência de um
\textit{volatility risk premium} positivo também nesse ambiente, ainda que com
características distintas das observadas em mercados mais maduros
\citep{hoang2023cryptoVRP}. Em particular, a elevada volatilidade intrínseca do
ativo, a negociação contínua e a presença recorrente de choques idiossincráticos
podem gerar maior instabilidade temporal e eventuais episódios de compressão ou
inversão do prêmio. Esta dissertação busca contribuir para esse campo ao oferecer
uma análise quantitativa abrangente do \textit{Bitcoin Volatility Risk Premium}
(BVRP), explorando suas propriedades estatísticas, sua variação sob diferentes
regimes de mercado e sua capacidade preditiva para retornos futuros por meio de
modelos modernos de aprendizado de máquina.

\section{Contextualização e motivação}

O mercado de opções de Bitcoin apresentou crescimento expressivo a partir de
2020, impulsionado pela maior participação institucional e pelo avanço da
infraestrutura de negociação. Dentre as plataformas existentes, a Deribit
consolidou-se como a principal bolsa de opções de Bitcoin em termos de volume e
\textit{open interest}, tornando-se referência para a construção de curvas de
volatilidade implícita e de índices derivados amplamente utilizados por
participantes de mercado.

A combinação entre elevada volatilidade intrínseca, negociação contínua (24/7) e
uma estrutura de liquidez ainda em consolidação torna o mercado de Bitcoin
particularmente propício ao estudo do VRP. Em princípio, espera-se que o prêmio
de risco de volatilidade seja mais elevado nesse ambiente, tanto pela maior
incerteza fundamental associada ao ativo quanto pela presença recorrente de
eventos idiossincráticos, como mudanças regulatórias, \textit{hacks},
liquidações de posições alavancadas e choques de liquidez em bolsas específicas.

Do ponto de vista da gestão de risco e da construção de estratégias
quantitativas, compreender a dinâmica do BVRP é relevante por, ao menos, três
razões principais:

\begin{enumerate}
    \item permite avaliar se o mercado precifica volatilidade de forma
    sistematicamente elevada, abrindo espaço para estratégias vendidas em
    volatilidade;
    
    \item fornece uma medida sintética da percepção de risco que pode ser
    incorporada em modelos de previsão de retornos e de \textit{asset
    allocation};

    \item auxilia na identificação de regimes de estresse em que o custo de
    proteção se torna especialmente elevado, contribuindo para decisões
    táticas e para o desenho de políticas eficazes de \textit{hedge}.
\end{enumerate}

\section{Problema de pesquisa}

À luz dessa contextualização, o problema de pesquisa que orienta esta
dissertação pode ser formulado nos seguintes termos:

\begin{quote}
    \textit{Quais são as principais propriedades estatísticas e dinâmicas do
    prêmio de risco de volatilidade em Bitcoin (BVRP) e em que medida essa
    variável contém informação útil para a previsão de retornos futuros do
    próprio Bitcoin em diferentes horizontes de tempo?}
\end{quote}

Essa questão desdobra-se em outras perguntas específicas, tais como:

\begin{itemize}
    \item O BVRP apresenta comportamento sistematicamente positivo ao longo da
          amostra analisada?
    \item Existem regimes bem definidos de alto e baixo BVRP associados a
          diferentes fases do ciclo de mercado?
    \item Em quais horizontes de tempo o BVRP parece apresentar maior poder
          preditivo para retornos do Bitcoin?
    \item Modelos de aprendizado de máquina conseguem explorar informações do
          BVRP, em conjunto com variáveis de estado de mercado, para melhorar a
          previsão de retornos em relação a abordagens lineares simples?
\end{itemize}

\section{Objetivos}

\subsection{Objetivo geral}

O objetivo geral desta dissertação é analisar, de forma quantitativa e
abrangente, o \textit{Bitcoin Volatility Risk Premium} (BVRP) e investigar sua
capacidade preditiva para retornos futuros do Bitcoin em horizontes de curto e
médio prazo, com ênfase na combinação entre métricas de volatilidade e modelos de
aprendizado de máquina.

\subsection{Objetivos específicos}

De forma mais detalhada, buscam-se os seguintes objetivos específicos:

\begin{enumerate}
    \item Construir uma série temporal consistente de volatilidade implícita a
          30 dias (IV30D) a partir de dados de opções de Bitcoin listadas na
          Deribit, bem como da volatilidade realizada a 30 dias (RV30D) a partir
          de retornos do mercado à vista.
    \item Estimar e analisar o BVRP (IV30D -- RV30D), caracterizando suas
          propriedades estatísticas e sua dinâmica temporal, com destaque para
          períodos de estresse de mercado.
    \item Investigar a relação entre o BVRP e retornos futuros do Bitcoin em
          diferentes horizontes temporais, por meio de análises gráficas,
          correlações e regressões lineares com erros-padrão robustos.
    \item Implementar modelos de aprendizado de máquina --- em particular LASSO,
          Ridge, \textit{Random Forest} e \textit{Gradient Boosting} --- com o
          objetivo de capturar potenciais não linearidades, interações complexas
          e efeitos de regime presentes na dinâmica do BVRP, avaliando ganhos
          preditivos fora da amostra em relação a modelos lineares de referência.
    \item Analisar o comportamento do BVRP sob diferentes regimes de volatilidade
          realizada, investigando se a distribuição do prêmio e o poder preditivo
          sobre retornos futuros variam sistematicamente entre ambientes de alta
          e baixa volatilidade.
    \item Discutir as implicações dos resultados para estratégias quantitativas,
          ressaltando oportunidades e riscos associados a operações que exploram
          o BVRP no mercado de opções de Bitcoin.
\end{enumerate}

\section{Hipóteses de trabalho}

Com base na literatura e na evidência preliminar, formulam-se as seguintes
hipóteses de trabalho:

\begin{description}
    \item[H1:] O BVRP é, em média, positivo ao longo da amostra analisada,
               refletindo um prêmio de proteção pago no mercado de opções de
               Bitcoin.
    \item[H2:] O BVRP apresenta regimes distintos de alto e baixo prêmio,
               associados a diferentes fases do ciclo de mercado e a episódios
               de aumento de incerteza.
    \item[H3:] O BVRP possui capacidade preditiva para retornos futuros do
               Bitcoin, especialmente em horizontes de médio prazo, ainda que o
               poder explicativo marginal seja moderado.
    \item[H4:] Modelos de aprendizado de máquina que incorporam o BVRP e
               variáveis de estado de mercado produzem previsões de retorno fora
               da amostra superiores às de regressões lineares que utilizam
               apenas variáveis tradicionais.
\end{description}

\section{Contribuições esperadas}

As contribuições deste trabalho podem ser agrupadas em três dimensões:

\begin{enumerate}
    \item \textbf{Contribuição empírica}: construção e análise detalhada de uma
          série de BVRP para Bitcoin em horizonte de 30 dias, utilizando dados de
          mercado de opções de alta liquidez e incorporando procedimentos de
          tratamento de dados específicos para o contexto de criptoativos.
    \item \textbf{Contribuição metodológica}: aplicação de modelos de aprendizado
          de máquina a um conjunto de variáveis que combina BVRP, medidas de
          volatilidade implícita e realizada e variáveis de estado de mercado,
          avaliando a utilidade incremental do BVRP em relação a abordagens
          lineares.
    \item \textbf{Contribuição aplicada}: discussão de implicações práticas para
          estratégias quantitativas e para a gestão de risco no mercado de
          opções de Bitcoin, incluindo limitações, riscos de cauda e desafios de
          implementação operacional.
\end{enumerate}

\section{Estrutura da dissertação}

Além deste capítulo introdutório, a dissertação está organizada da seguinte forma:

\begin{itemize}
    \item O Capítulo~2 apresenta o referencial teórico que fundamenta o estudo,
          discutindo conceitos de volatilidade implícita e realizada, prêmio de
          risco de volatilidade, especificidades do mercado de opções de
          criptoativos e aplicações recentes de modelos de aprendizado de
          máquina em finanças.

    \item O Capítulo~3 descreve as bases de dados utilizadas e apresenta uma
          análise exploratória das principais variáveis de interesse, incluindo
          a construção das séries de volatilidade implícita (IV30D) e realizada
          (RV30D) e as estatísticas descritivas do BVRP.

    \item O Capítulo~4 apresenta a metodologia empírica adotada, especificando
          os modelos econométricos, o protocolo de validação fora da amostra e
          as variáveis de estado de mercado utilizadas ao longo da dissertação.

    \item O Capítulo~5 apresenta os primeiros resultados econométricos, obtidos
          por meio de regressões OLS com erros-padrão de Newey--West, avaliando
          o conteúdo informacional do BVRP para a previsão de retornos futuros
          do Bitcoin em múltiplos horizontes.

    \item O Capítulo~6 investiga a previsibilidade do BVRP fora da amostra por
          meio de modelos de \textit{machine learning} --- LASSO, Ridge,
          \textit{Random Forest} e \textit{Gradient Boosting} ---, comparando
          seu desempenho a \textit{benchmarks} baseados na média histórica do
          prêmio.

    \item O Capítulo~7 analisa a utilidade econômica do BVRP como sinal para
          estratégias condicionais de negociação no mercado à vista, avaliando
          estratégias baseadas em quantis e em regimes de volatilidade realizada
          por meio de métricas de desempenho ajustadas ao risco.

    \item O Capítulo~8 retoma a análise preditiva do BVRP em uma abordagem
          supervisionada mais estruturada, motivada pelas limitações das regras
          determinísticas do capítulo anterior. Compara-se o desempenho de quatro
          classes de modelos supervisionados sob um desenho experimental rigoroso,
          com ênfase na importância das variáveis preditoras e na estabilidade
          temporal das estimativas.

    \item O Capítulo~9 analisa o comportamento do BVRP sob diferentes regimes de
          volatilidade realizada, caracterizando a distribuição do prêmio por
          regime, testando diferenças nos retornos futuros e investigando o poder
          preditivo incremental da interação entre o BVRP e o ambiente de
          volatilidade.

    \item Por fim, as Considerações Finais sintetizam as principais evidências
          encontradas ao longo do trabalho, discutem as limitações da análise e
          apontam extensões para pesquisas futuras.
\end{itemize}
